\section{Introduction}
\label{sec:intro}

Retrieval-augmented generation has become the default way to put a language model
in front of a document corpus \citep{lewis2020rag,gao2023ragsurvey}. The standard
recipe embeds passages, retrieves the nearest neighbours of an embedded question,
and lets the model answer from what came back. It works well when the answer is
contained in a passage that resembles the question.

A large class of enterprise questions is not like that. Consider a question a
financial analyst would actually ask: \emph{for every company where a particular
person is an officer or director, list that company's disclosed subsidiaries and
summarize the cybersecurity risks it reported.} Three facts are required. The
person's role is disclosed in an insider-ownership filing. The subsidiaries are
disclosed in an exhibit attached to the annual report. The risks are narrative
prose inside a numbered section of that same report. These are three separate
documents, in three formats, produced by three different processes.

The difficulty is not that the passages are hard to rank. It is that no passage
contains the combination. We measured this directly on our corpus: the surname of
the subject of our hardest question appears in \textbf{0 of 5{,}210} extracted
risk-factor sections, and only \textbf{1 of 151} subsidiary names appears anywhere
in its own parent's risk-factor section. Ownership and risk are, empirically,
disjoint. \Cref{fig:containers} shows the structure that produces this.

This has a consequence that is easy to state and easy to overlook: a better
encoder cannot fix it, and neither can a longer context window. Retrieval can
only return text that exists. If the conjunction of facts a question needs is
never co-located in any passage, then ranking passages more accurately retrieves
the same insufficient evidence. The ceiling is a property of the corpus, not of
the retriever.

There is a second and subtler failure. Similarity search has no mechanism for
filtering by entity, because ``belongs to company $X$'' is not a semantic property
of a passage. Two risk-factor sections about supply-chain disruption are near each
other in embedding space whether or not they belong to the company being asked
about. At small scale this rarely bites, because few candidates compete. We show
that it bites hard at 508{,}714 chunks drawn from 5{,}210 filers, and --- more
usefully --- that it bites on exactly the questions where similarity search is
supposed to be the right tool.

The obvious alternative is to store the structure explicitly: resolve entities to
stable identifiers, materialize the relationships as edges, and attach a
provenance key to every edge so that any claim can be traced to the document that
supports it. This is the family of designs now loosely called GraphRAG
\citep{edge2024graphrag,peng2024graphragsurvey}. The question this paper addresses
is not whether that is possible but whether it is \emph{worth it}, and on which
question shapes --- measured against two baselines that were given every
reasonable advantage.

We are deliberately narrow. This is an empirical systems study on one corpus with
one model and twenty questions. It contributes measurements and diagnosed failure
modes, not a new algorithm. Our contributions are:

\begin{itemize}
  \item \textbf{A corpus and benchmark with hard entity identity.} Twelve monthly
        EDGAR feeds ingested into a relational store and a graph
        (\Cref{tab:corpus}), with 20 questions in four categories. The multi-hop
        subjects were selected on measured evidence: the surname must appear in
        \emph{zero} risk-factor sections and must be shared by two or more
        distinct filer identifiers, so that resolving identity by name is
        demonstrably wrong rather than merely inelegant
        (\Cref{sec:design}).
  \item \textbf{A matched three-way comparison.} 180 programmatically scored
        attempts under one model at temperature 0, with the fan-out guardrail and
        repair policy applied identically to both query-writing arms
        (\Cref{sec:arch}), and with two of the four categories chosen so that the
        graph should lose.
  \item \textbf{Two results that contradict our expectation.} Vector RAG lost the
        category it was expected to win, and plain text-to-SQL beat GraphRAG on
        cross-document entity resolution. We report both with the same prominence
        as the wins and diagnose each (\Cref{sec:failures}).
  \item \textbf{A cost accounting.} GraphRAG's accuracy is bought:
        780{,}500 tokens against text-to-SQL's 158{,}171, and a repair call
        triggered on 34 of 60 attempts against 7 of 60
        (\Cref{tab:main}, \Cref{tab:repair}). We also report the reverse cost ---
        text-to-SQL's generated joins averaged 92\,s and peaked at 445\,s on the
        multi-hop category.
  \item \textbf{Eight defects found in the measurement harness itself.} Each
        would have produced a confident and wrong conclusion had it gone
        unnoticed, including one ground-truth error found after the results were
        first computed (\Cref{sec:repro}). We argue that an unaudited benchmark
        harness is itself a leading source of false findings.
\end{itemize}

\Cref{sec:related} places the work; \Cref{sec:dataset} and \Cref{sec:arch}
describe the corpus and the three architectures; \Cref{sec:design} gives the
protocol; \Cref{sec:results} reports what happened; \Cref{sec:failures} explains
the failures; \Cref{sec:threats} states what these measurements cannot support.
